\chapter{Fundamentos}
\label{cap:fundamentos}

\section{El motor: CFM56-7B26}
\label{sec:motor}

El CFM56-7B (CFM International, consorcio Safran--GE) propulsa la familia Boeing 737NG y es uno
de los motores comerciales más extendidos de la historia, lo que garantiza abundancia de datos
públicos de tipo. Se elige la variante \textbf{-7B26} (\SI{117}{\kilo\newton} de empuje de
despegue, Boeing 737-800) por ser la más común. Según su hoja de datos de certificado de tipo
(TCDS EASA E.004, ed.~07~\cite{EASA2023}), es un turbofan de doble rotor y alta relación de
derivación: fan de una etapa, compresor de baja (\emph{booster}) de 3 etapas, compresor de alta
(HPC) de 9 etapas, cámara de combustión anular, turbina de alta (HPT) de 1 etapa y turbina de
baja (LPT) de 4 etapas, con control digital FADEC de doble canal. El parámetro primario de
empuje es N1 (velocidad del rotor de baja), no EPR.

\begin{table}[htbp]
  \centering\small
  \caption{Datos certificados del CFM56-7B26 empleados como contrato de calibración. Fuentes:
  TCDS EASA E.004 ed.~07~\cite{EASA2023} y banco de datos de emisiones OACI, ed.~03/2026, UID
  3CM033~\cite{ICAO2026}.}
  \label{tab:tcds}
  \begin{tabular}{llll}
    \toprule
    Magnitud & Valor & Fuente & Estado \\
    \midrule
    Empuje de despegue & \SI{11699}{\deca\newton} (= \SI{116.99}{\kilo\newton}) & TCDS & verificado \\
    Empuje máximo continuo & \SI{11521}{\deca\newton} & TCDS & verificado \\
    N1 límite (104\,\%) & \SI{5382}{rpm} $\Rightarrow$ 100\,\% = \SI{5175}{rpm} & TCDS & verificado \\
    N2 límite (105\,\%) & \SI{15183}{rpm} $\Rightarrow$ 100\,\% = \SI{14460}{rpm} & TCDS & verificado \\
    EGT máx.\ despegue / continuo & \SI{950}{\degreeCelsius} / \SI{925}{\degreeCelsius} (mostrada) & TCDS & verificado \\
    Peso en seco & \SI{2386}{\kilo\gram} & TCDS & verificado \\
    Relación de derivación & 5{,}1 & EEDB & verificado \\
    Relación de presiones (SLS, nominal) & 27{,}61 & EEDB & verificado \\
    $W_f$ a 100/85/30/7\,\% $F_{00}$ & 1{,}221 / 0{,}999 / 0{,}338 / 0{,}113 \si{\kilo\gram\per\second} & EEDB & verificado \\
    \bottomrule
  \end{tabular}
\end{table}

Dos matices del TCDS resultan relevantes para el modelado y suelen pasarse por alto:
\begin{itemize}
  \item La EGT certificada se mide en la estación \textbf{T49.5} (tobera de la segunda etapa de
    la LPT), no en la salida de la HPT; y el valor \emph{mostrado} en cabina añade una función
    \emph{shunt} de $+\SI{30}{\degreeCelsius}$ (motores SAC, activa por encima de \SI{8500}{rpm} de
    N2) más un \emph{trim} dependiente del modelo (nulo para el -7B26).
  \item Los límites de sangrado están tabulados por estación de extracción (5.ª y 9.ª etapa del
    HPC) y régimen; a régimen nominal el sangrado de 9.ª etapa se limita al 7\,\% del flujo
    primario.
\end{itemize}

\section{Estaciones y parámetros corregidos}
\label{sec:estaciones}

Se adopta la numeración SAE ARP755A~\cite{SAE755}: 2 entrada del fan; 21 salida del fan; 25
entrada del HPC; 3 salida del HPC; 4 salida de cámara; 45 salida de la HPT; 5 salida de la LPT.
El rotor de baja (N1) une fan, \emph{booster} y LPT; el de alta (N2), HPC y HPT.

Para comparar vuelos en condiciones distintas, las medidas se corrigen a condiciones ISA a nivel
del mar con $\theta = T_{t2}/\SI{288.15}{\kelvin}$ y $\delta = P_{t2}/\SI{101.325}{\kilo\pascal}$:
\begin{equation}
  N_{\mathrm{corr}} = \frac{N}{\sqrt{\theta}}, \qquad
  W_{f,\mathrm{corr}} = \frac{W_f}{\delta\,\theta^{a}}, \qquad
  \mathrm{EGT}_{\mathrm{corr}} = \frac{\mathrm{EGT}}{\theta},
\end{equation}
donde el exponente $a$ (habitualmente $0{,}5$--$0{,}7$) no se toma de tablas genéricas sino que
se ajustará por regresión sobre barridos del propio modelo termodinámico, eliminando una fuente
clásica de error sistemático en las líneas base.

\section{Formulación del \emph{Gas Path Analysis}}
\label{sec:gpa}

El estado de salud del camino de gas se describe con el vector de desviaciones respecto a motor
nuevo
\begin{equation}
  \mathbf{x} = \left[\Delta\eta_{\mathrm{fan}}, \Delta\Gamma_{\mathrm{fan}},
  \Delta\eta_{\mathrm{LPC}}, \Delta\Gamma_{\mathrm{LPC}},
  \Delta\eta_{\mathrm{HPC}}, \Delta\Gamma_{\mathrm{HPC}},
  \Delta\eta_{\mathrm{HPT}}, \Delta\Gamma_{\mathrm{HPT}},
  \Delta\eta_{\mathrm{LPT}}, \Delta\Gamma_{\mathrm{LPT}}\right]^{\mathsf T},
\end{equation}
donde $\Delta\eta$ es la desviación relativa de eficiencia isentrópica y $\Delta\Gamma$ la de
capacidad de flujo corregida de cada turbomáquina. Las medidas disponibles se expresan como
desviaciones $\Delta\mathbf{z}$ de los parámetros corregidos respecto a la línea base en el
mismo punto de operación $\mathbf{u}$ (altitud, Mach, N1, $\Delta T_{\mathrm{ISA}}$, sangrado).
El modelo lineal local es
\begin{equation}
  \Delta\mathbf{z} = \mathbf{H}(\mathbf{u})\,\mathbf{x} + \mathbf{S}\,\mathbf{b} + \mathbf{v},
  \qquad \mathbf{v}\sim\mathcal{N}(\mathbf{0},\mathbf{R}),
  \label{eq:gpa-lineal}
\end{equation}
con $\mathbf{H}$ la \emph{matriz de coeficientes de influencia} (ICM), jacobiano
$\partial\mathbf{z}/\partial\mathbf{x}$ evaluado en $\mathbf{u}$ y obtenido por diferencias
centrales sobre el modelo termodinámico; $\mathbf{b}$ los sesgos de sensor y $\mathbf{R}$ la
covarianza del ruido de medida.

\paragraph{Estimación WLS.} Con 4 medidas de cabina (N1, N2, EGT, $W_f$) y 10 incógnitas el
sistema es infra-determinado ($\operatorname{rango}(\mathbf{H})\leq 4$), por lo que la solución
de mínimos cuadrados ponderados exige regularización con una prior $\mathbf{P}_0$ de magnitudes
de deterioro esperables:
\begin{equation}
  \hat{\mathbf{x}} = \left(\mathbf{H}^{\mathsf T}\mathbf{R}^{-1}\mathbf{H}
  + \lambda\,\mathbf{P}_0^{-1}\right)^{-1}\mathbf{H}^{\mathsf T}\mathbf{R}^{-1}\Delta\mathbf{z}.
  \label{eq:wls}
\end{equation}

\paragraph{Seguimiento con filtro de Kalman.} El deterioro evoluciona lentamente, lo que
justifica un modelo de paseo aleatorio:
\begin{equation}
  \mathbf{x}_k = \mathbf{x}_{k-1} + \mathbf{w}_k, \quad \mathbf{w}\sim\mathcal{N}(\mathbf{0},\mathbf{Q});
  \qquad
  \mathbf{z}_k = \mathbf{H}(\mathbf{u}_k)\,\mathbf{x}_k + \mathbf{v}_k,
  \label{eq:kalman}
\end{equation}
con $\mathbf{Q}$ calibrada a las tasas de degradación y estado aumentable con sesgos de sensor
\cite{Volponi2014}. En los eventos registrados de lavado de compresor se reinicia la parte
recuperable del estado.

\paragraph{\emph{Smearing} y observabilidad.} Cuando $\mathbf{H}$ está mal condicionada, el
ruido proyecta un fallo concentrado en un solo componente sobre varios: es el fenómeno de
\emph{smearing}, debilidad clásica del GPA lineal~\cite{Li2002}. El análisis SVD de la ICM
(rango, número de condición, ángulos entre firmas de fallo) cuantifica esta debilidad y define
el subconjunto ``confusable'' usado por la hipótesis H2: pares de fallos cuyas firmas forman un
ángulo inferior a $15^{\circ}$.

\section{Física de la degradación}
\label{sec:degradacion}

\begin{table}[htbp]
  \centering\small
  \caption{Mecanismos de degradación modelados y efectos típicos según la
  literatura~\cite{Diakunchak1992,Kurz2012}.}
  \label{tab:degradacion}
  \begin{tabular}{p{3.2cm}p{2.2cm}p{3.6cm}p{2.6cm}p{1.6cm}}
    \toprule
    Mecanismo & Componente & Efecto típico & Perfil temporal & Recuperable \\
    \midrule
    \emph{Fouling} & fan/LPC/HPC & $\Delta\Gamma$ $-1$\ldots$-4$\,\%; $\Delta\eta$ $-0{,}5$\ldots$-2{,}5$\,\% & exponencial saturante & lavado (30--70\,\%) \\
    Erosión & compresores & $\Delta\eta$ $-0{,}5$\ldots$-1{,}5$\,\% & lineal lento & no \\
    Holgura de punta & HPC/HPT & $\Delta\eta$ $-0{,}5$\ldots$-1{,}5$\,\% & rápido inicial, luego lento & no \\
    Deterioro sección caliente & HPT & $\Delta\eta$ $-1$\ldots$-3$\,\%; $\Delta\Gamma$ $+1$\ldots$+2$\,\% & lineal/acelerado & no \\
    FOD & fan/HPC & escalón $\Delta\eta$ $-0{,}5$\ldots$-2$\,\% & escalón (Poisson) & no \\
    Deriva de sensor & cualquier sonda & p.\ ej.\ EGT $+1$\ldots$+5$~\si{\degreeCelsius}/1000 ciclos & rampa & recalibración \\
    \bottomrule
  \end{tabular}
\end{table}

El efecto agregado observable es la erosión del \emph{margen de EGT} (distancia entre la EGT en
despegue de día caliente y su límite certificado), desde \SI{70}{}--\SI{100}{\degreeCelsius} en motor
nuevo hasta cero al final de la vida útil en ala, con el patrón característico ``rápido al
principio, lento después, dientes de sierra por lavados''. La \cref{tab:degradacion} resume los
mecanismos, sus firmas y sus perfiles temporales; sus magnitudes parametrizan el generador de
flota sintética de la fase F2.
